\section{Protobuf Message Design}\label{sec:proto-msg}

Mirroring the descriptors are the values, which carry actual data rather than
describe the shape of a message.

\subsection{Values in Lean}

In the Lean formalization, the structure of a value matches the structure of
descriptors across all three levels.

\begin{itemize}
\item A descriptor can be matched to a value.
\item A field can be matched to a slot.
\item A field type can be matched to a payload.
\end{itemize}

\begin{minted}{lean}
mutual
inductive Value   where | mk (es : List ((_ : Int) × Val))
inductive Slot    where | implicit (p : Payload)
                        | optional (p : Option Payload)
                        | repeated (ps : List Payload)
inductive Payload where | int (z : Int) | bool (b : Bool) | string (s : String)
                        | bytes (bs : List UInt8) | float (bits : UInt32)
                        | double (bits : UInt64) | msg (v : Value)
end
\end{minted}

Notice that just like with the descriptor, the presence wrapper sites above the
payload, and also reduces the number of constructor cases (marginally) and
creates a regular structure for \texttt{optional} and \texttt{repeated} fields
of an \texttt{option} / \texttt{list} of payloads rather than an \texttt{option}
/ \texttt{list} with a specific instance type. Actual data is represented in
Lean's native types. The \texttt{msg} payload constructor is the only major
structural difference from the descriptor, since it sits in the \texttt{Payload}
type rather than in a separate mutually inductive type. (I've eliminated some of
the difference in the formal grammar, but the distinction on the lean side is
worth noting.) However, the inductive structure \emph{does} match --- it's
under the presence modifier, but still above the leaf of the descriptor or value
--- and the descriptor's scalar type gains derivable equality by being outside the
mutually recursive block. The way that the value \texttt{Payload}'s are used,
this isn't required, they're compared against a field from the descriptor in a
way to avoid this problem. Also, doubles and floats are represented as unsigned
integers to get bit-level equality. Lean's formalization of IEEE 754 floating
point numbers is opaque and the defined equality isn't reflexive (NaN $\ne$ NaN). 

{\color{red} This might change. It's possible to pull the scalar payloads out if
needed.}

\begin{figure}[h]
  \centering
  \begin{syntax}
    \categoryFromSet[Integers]{z}{\mathbb{Z}}

    \categoryFromSet[Booleans]{b}{\mathbb{B}}

    \abstractCategory[Byte Strings]{bs}

    \abstractCategory[Strings]{\mathrm{s}}

    \categoryFromSet[32-bit Unsigned Integers]{\mathrm{u^{32}}}{[0, 2^{32})}

    \categoryFromSet[64-bit Unsigned Integers]{\mathrm{u^{64}}}{[0, 2^{64})}

    \category[Payload]{p}
    \alternative{z}
    \alternative{b}
    \alternative{\mathrm{s}}
    \alternative{bs}
    \alternative{\mathrm{u^{32}}}
    \alternative{\mathrm{u^{64}}}
    \alternative{v}

    \category[Slot]{s}
    \alternative{\impt{p}}
    \alternative{\optt{\text{none}}}
    \alternative{\optt{p}}
    \alternative{\rept{[p_1; \dots; p_m]}}

    \category[Value List]{l}
    \alternative{\texttt{[]}}
    \alternative{(n, s) :: l}

    \category[Value]{v}
    \alternative{\llangle l \rrangle_v}
  \end{syntax}
  \caption{Pollux syntax for Protobuf values.}\label{fig:proto-val-syn}
\end{figure}

\subsection{Value Totality}\label{sec:value-totality}

In real-world Protobuf implementations, a value always has the exact same domain
as the descriptor which generated it. If a message has a missing field (for any
reason) and that field is marked as implicit in the descriptor, the default
value is automatically populated into the field. Pollux matches this model by
requiring that when a value is linked to a specific descriptor, the domains of
the two match exactly. This will be part of the value well-formedness
conditions. Pollux will model a missing field using either the default value for
that type, \texttt{none} for optional values and an empty list for
\texttt{repeated} ones, just like Protobuf.

This choice strengthens the connection between a descriptor and values under it,
since rather then being able to mix and match the two arbitrarily --- something
that no real-world Protobuf implementation allows --- the assumptions made by
Pollux closely model the data assumptions of Protobuf.

However, on the proof aspect of Pollux both directions are important and worth
concerning.
\begin{itemize}
\item \textbf{Declared implies present.} A field being declared in the
  descriptor implies that something is present here in the matching value. This
  is the \emph{no omissions requirement}.
\item \textbf{Present implies declared.} Likewise a value being present in the
  message implies that it was declared in the descriptor. This is the \emph{no
    unknown values requirement}.
\end{itemize}

Starting with the first case, this requirement arises from some trouble with the
intermediate parsing format's compatibility relation, specifically the
\textsc{M-Declare} rule, where the writer's descriptor carries a field that
isn't present in the value, and the reader gets a missing value (a
intermediate format feature not present in Protobuf, this would become the
default for the type of the field).
\[
  \infer[M-Declare]{m_1 : d_1 \preceq m_2 : d_2 \\\\ k \not \in \dom{d_1} \\ k \not \in
    \dom{m_1} \\ f_1 \propto f_2}{ m_1 : d_1[k \mapsto f_1] \preceq m_2[k \mapsto \mathtt{M\_MISSING}] :
    d_2[k \mapsto f_2]}
\]
For that format, this rule was required since without it, the message
compatibility relation could not follow the ID compatibility relation which
characterizes what could happen to a value being serialized and parsed under the
same descriptor. The well-formedness condition on the values in the intermediate
format only constrained the value without regard for a descriptor.

Under the new totality requirement, this rule is vacuous. But why is removing it
a good thing?
\begin{enumerate}
\item This doesn't happen under the Protobuf data model, which would have
  inserted the default value when creating or parsing the value originally. This
  rule only exists since the well-formedness condition of the intermediate
  format was too weak, but becomes a case in every proof about the message
  compatibility relation that describes a case that Protobuf will never express.
\item Without totality, the intermediate format never had message equality when
  serializing and parsing under the same descriptor, and this requirement
  simplifies the proofs and theorem statements.
\end{enumerate}

In the second case, the no unknown fields requirement does two important
things:

\begin{enumerate}
\item The transformations output (Section~\ref{sec:proto-transform}) will always
  have a domain of $\dom(d_2)$ since we want to provide a value when viewed
  against $d_2$. So the keys that we lose when transforming a value $v$ are
  exactly $\dom{v} \setminus \dom{d_2}$. Having exact domain equality between
  $v$ and $d_1$ means that this is the same as $\dom{d_1} \setminus \dom{d_2}$
  and enables the descriptor level compatibility relation to characterize the
  value-level transformation. The drop rule in both the descriptor compatibility
  relation and the message one for Protobuf will be applied only at write-only
  descriptor key. Allowing unknown fields to exist in the value would break
  this, and a descriptor analysis wouldn't be able to meaningful reason about
  what happens at the value level.
\item Transforming a value from it's linked descriptor to that same descriptor
  should yield exactly the same value. Without this requirement, an unknown field
  would be dropped since Pollux doesn't track unknown fields explicitly
  (Protobuf actually does, for largely this reason, motivated by a potential
  man-in-the-middle where the sender and receiver both have a larger descriptor
  then the relay).
\end{enumerate}

At the byte level, the serializer will never emit a tag that the parser using
the same descriptor can't understand, so our prospective is that an unknown
field must come from a schema difference in the descriptors (or some
modification made to the value on the network; well outside our scope).

\subsection{Cardinality and Scalar Carriers}

The presence marker in the value layer has three constructors, while the
cardinality marker in the descriptor layer has four, with \texttt{oneof} being
the missing one. Be default in Protobuf, all \texttt{oneof} fields are treated
as \texttt{optional}, which Pollux also respects. The requirement that only one
is set is a cross-field constraint, which will live in value's legal predicate.

Additionally, while the descriptor scalar types enumerate all Protobuf type, the
value collapses several of these into a shared carrier. Specifically for all the
various integer types, the value uses a mathematical integer when the Protobuf
scalar types change in encoding and range. The legal predicate for values will
enforce the range requirement while the encoding choice is only relevant
during the serialization and parsing process, not the final values.

\subsection{Value Initialization}

Since messages have to match their descriptor, it's important that we're able to
generate initial, empty messages to establish this invariant. Notice that the
\texttt{Field.init} function takes the cardinality information in addition to the
type information. This is what lets the initialization process differentiate
between the default value of a repeated field (an empty list), optional field
(\texttt{none}) and scalar field.

\begin{minted}{lean}
def Field.init : Field → Val
  | .mk .singular (.scalar s) => .implicit (ScalarType.defaultPayload s)
  | .mk .singular (.msg _)    => .optional none   -- unreachable when Legal
  | .mk .optional _           => .optional none
  | .mk .repeated _           => .repeated []
  | .mk (.oneof _) _          => .optional none

def Value.init (d : Desc) : Value := ⟨d.entries.map
    (fun e => ⟨e.1, Field.init e.2⟩)⟩

theorem Value.get?_init (d : Desc) (k : Int) :
    (Value.init d).get? k = (d.get? k).map Field.init
\end{minted}

This is one of two places where the descriptor seal is broken and (recursive) iteration over
the descriptor is performed, the other being the injection of missing fields
during parsing. 

\subsection{Value Validity} 

Descriptor validity broke down into two requirements, a Lean specific invariant
on the internal structure of our representation, and the Legal predicate that
enforces some particular requirements of the Protobuf compiler. A value being
well-formed is a bit more complex since we require that it's well-formed
\emph{against} a descriptor. The value well-formedness predicate will be used as
the output well-formedness predicate of the Protobuf serializer, meaning we'll
be able to assume that a value input to the serializer is well-formed and
use that in combination progress theorem that all well-formed values will
serialize. Validity on the value is the precondition for the top-level
round-trip theorems, and are encoded into the type of the serializer from our
serializer combinator library. The actual contents of the validity predicate are
stronger than needed to show that the serialization will always succeed, but can
be weakened to imply the needed combinator well-formedness conditions throughout the
combinator library.

\vspace{1em}
\begin{definition}[Valid Message]\label{def:vvalid}
  A message $m$ is valid against a descriptor $d$, denoted $v \checkmark^d$ if
  both $v \checkmark_w$ and $v \checkmark_l^d$. 
\end{definition}

\vspace{1em}
\begin{definition}[Message Well-Formedness]\label{def:vwf}
  A message $m$ is well-formed, denoted $m \checkmark_w$, if all fields in
  the message have a unique field number, the fields are stored in sorted
  order, and all nested messages are also well-formed.
\end{definition}

\vspace{1em}
\begin{definition}[Legal Message]\label{def:vlegal}
  A message $m$ is legal with regards to descriptor $d$, denoted $m
  \checkmark_l^d$, if
  \begin{itemize}
  \item Totality, $\dom{m} = \dom{d}$.
  \item For all values in the message
    \begin{itemize}
    \item The presence of the field agrees with the cardinality in the
      descriptor.
    \item The type of the value agrees with the type of the field in the
      descriptor. For integer scalars, this condition also enforces the range
      bounds associated with the various Protobuf integer types (i.e.\ \ints{},
      \uintl{}, etc.) while for nested messages, that means $v' \checkmark^{d'}$.
    \end{itemize}
  \item At most one key in a \texttt{oneof} group defined in the descriptor is
    set. This is the only non-pairwise requirement that a well-formed message has
    to respect.
  \end{itemize}
\end{definition}

Descriptor legality (Definition~\ref{def:dlegal}), specifically the fact that
messages have to be \texttt{optional}, is important here since it makes it
possible to initialize a recursive message, and we can assume that anytime a
message is present, it's a non-default. This means that the Pollux check for if
a value is a default can be decidable, even though descriptors and fields don't
have a decidable equality instance defined.

\vspace{1em}
\begin{lemma}
  For all descriptors $d$, $\mathrm{init}(d) \checkmark^d$.
\end{lemma}

\vspace{1em}
\begin{definition}[Message Size]\label{def:vsize}
  The size of a message, $\mathrm{size}_v$, is defined below. 
  \begin{mathpar}
    \mathrm{size}_v(\llangle l \rrangle) = 1 + \mathrm{size}_v(l)

    \mathrm{size}_v(l) =
    \begin{cases}
      0 & \text{if } l = \texttt{[]} \\
      \mathrm{size}_v(v) + \mathrm{size}_v(l') & \text{if } l = (\_, v) :: l'
    \end{cases}

    \mathrm{size}_v(v) =
    \begin{cases}
      1 + \mathrm{size}_v(p) & \text{if } v = \impt{p} \\
      1 & \text{if } v = \optt{\text{none}} \\
      1 + \mathrm{size}_v(p) & \text{if } v = \optt{p} \\
      1 + \mathrm{size}_v(p_1) + \dots + \mathrm{size}_v(p_n) & \text{if } v = \rept{[p_1;\dots;p_n]}
    \end{cases}

    \mathrm{size}_v(p) =
    \begin{cases}
      \mathrm{size}_v(\llangle l' \rrangle) & \text{if } p = \llangle l'
                                              \rrangle \\
      1 & \text{otherwise}
    \end{cases}
  \end{mathpar}
\end{definition}

% Local Variables:
% citar-bibliography: ("../pollux.bib")
% TeX-master: "../pollux.tex"
% TeX-command-default: "Make"
% jinx-local-words: "protobuf"
% End:
